%**************************************%
%*    Generated from PreTeXt source   *%
%*    on 2020-04-09T18:59:59+05:30    *%
%*                                    *%
%*      https://pretextbook.org       *%
%*                                    *%
%**************************************%
\documentclass[oneside,10pt,]{book}
%% Custom Preamble Entries, early (use latex.preamble.early)
%% Default LaTeX packages
%%   1.  always employed (or nearly so) for some purpose, or
%%   2.  a stylewriter may assume their presence
\usepackage{geometry}
%% Some aspects of the preamble are conditional,
%% the LaTeX engine is one such determinant
\usepackage{ifthen}
%% etoolbox has a variety of modern conveniences
\usepackage{etoolbox}
\usepackage{ifxetex,ifluatex}
%% Raster graphics inclusion
\usepackage{graphicx}
%% Color support, xcolor package
%% Always loaded, for: add/delete text, author tools
%% Here, since tcolorbox loads tikz, and tikz loads xcolor
\PassOptionsToPackage{usenames,dvipsnames,svgnames,table}{xcolor}
\usepackage{xcolor}
%% begin: defined colors, via xcolor package, for styling
%% end: defined colors, via xcolor package, for styling
%% Colored boxes, and much more, though mostly styling
%% skins library provides "enhanced" skin, employing tikzpicture
%% boxes may be configured as "breakable" or "unbreakable"
%% "raster" controls grids of boxes, aka side-by-side
\usepackage{tcolorbox}
\tcbuselibrary{skins}
\tcbuselibrary{breakable}
\tcbuselibrary{raster}
%% We load some "stock" tcolorbox styles that we use a lot
%% Placement here is provisional, there will be some color work also
%% First, black on white, no border, transparent, but no assumption about titles
\tcbset{ bwminimalstyle/.style={size=minimal, boxrule=-0.3pt, frame empty,
colback=white, colbacktitle=white, coltitle=black, opacityfill=0.0} }
%% Second, bold title, run-in to text/paragraph/heading
%% Space afterwards will be controlled by environment,
%% independent of constructions of the tcb title
%% Places \blocktitlefont onto many block titles
\tcbset{ runintitlestyle/.style={fonttitle=\blocktitlefont\upshape\bfseries, attach title to upper} }
%% Spacing prior to each exercise, anywhere
\tcbset{ exercisespacingstyle/.style={before skip={1.5ex plus 0.5ex}} }
%% Spacing prior to each block
\tcbset{ blockspacingstyle/.style={before skip={2.0ex plus 0.5ex}} }
%% xparse allows the construction of more robust commands,
%% this is a necessity for isolating styling and behavior
%% The tcolorbox library of the same name loads the base library
\tcbuselibrary{xparse}
%% Hyperref should be here, but likes to be loaded late
%%
%% Inline math delimiters, \(, \), need to be robust
%% 2016-01-31:  latexrelease.sty  supersedes  fixltx2e.sty
%% If  latexrelease.sty  exists, bugfix is in kernel
%% If not, bugfix is in  fixltx2e.sty
%% See:  https://tug.org/TUGboat/tb36-3/tb114ltnews22.pdf
%% and read "Fewer fragile commands" in distribution's  latexchanges.pdf
\IfFileExists{latexrelease.sty}{}{\usepackage{fixltx2e}}
%% Text height identically 9 inches, text width varies on point size
%% See Bringhurst 2.1.1 on measure for recommendations
%% 75 characters per line (count spaces, punctuation) is target
%% which is the upper limit of Bringhurst's recommendations
\geometry{letterpaper,total={340pt,9.0in}}
%% Custom Page Layout Adjustments (use latex.geometry)
%% This LaTeX file may be compiled with pdflatex, xelatex, or lualatex executables
%% LuaTeX is not explicitly supported, but we do accept additions from knowledgeable users
%% The conditional below provides  pdflatex  specific configuration last
%% The following provides engine-specific capabilities
%% Generally, xelatex is necessary for non-Western fonts
\ifthenelse{\boolean{xetex} \or \boolean{luatex}}{%
%% begin: xelatex and lualatex-specific configuration
\ifxetex\usepackage{xltxtra}\fi
%% realscripts is the only part of xltxtra relevant to lualatex 
\ifluatex\usepackage{realscripts}\fi
%% fontspec package provides extensive control of system fonts,
%% meaning *.otf (OpenType), and apparently *.ttf (TrueType)
%% that live *outside* your TeX/MF tree, and are controlled by your *system*
%% (it is possible that a TeX distribution will place fonts in a system location)
\usepackage{fontspec}
%% We use Latin Modern (lmodern) as the default font
%% So we check that it is available as a system font
\IfFontExistsTF{Latin Modern Roman}{}{\GenericError{}{The font "Latin Modern Roman" requested by PreTeXt output is not available as a system font}{Consult the PreTeXt Guide for help with LaTeX fonts.}{}}
%% We then define various font family commands using a vanilla version,
%% with the intention of letting a style override these choices
%% \setmainfont can be re-issued, and \renewfontfamily can redefine others
\setmainfont{Latin Modern Roman}[SmallCapsFont={Latin Modern Roman Caps}, SlantedFont={Latin Modern Roman Slanted}]
\newfontfamily{\divisionfont}{Latin Modern Roman}
\newfontfamily{\blocktitlefont}{Latin Modern Roman}
\newfontfamily{\contentsfont}{Latin Modern Roman}
\newfontfamily{\pagefont}{Latin Modern Roman}[SlantedFont={Latin Modern Roman Slanted}]
\newfontfamily{\tabularfont}{Latin Modern Roman}[SmallCapsFont={Latin Modern Roman Caps}]
\newfontfamily{\xreffont}{Latin Modern Roman}
\newfontfamily{\titlepagefont}{Latin Modern Roman}
%% begin: font information supplied by "font-xelatex-style" template
%% end: font information supplied by "font-xelatex-style" template
%% 
%% Extensive support for other languages
\usepackage{polyglossia}
%% Set main/default language based on pretext/@xml:lang value
%% document language code is "en-US", US English
%% usmax variant has extra hypenation
\setmainlanguage[variant=usmax]{english}
%% Enable secondary languages based on discovery of @xml:lang values
%% Enable fonts/scripts based on discovery of @xml:lang values
%% Western languages should be ably covered by Latin Modern Roman
%% end: xelatex and lualatex-specific configuration
}{%
%% begin: pdflatex-specific configuration
\usepackage[utf8]{inputenc}
%% PreTeXt will create a UTF-8 encoded file
%% begin: font setup and configuration for use with pdflatex
%% Portions of a document, are, or may, be affected by font-changing commands
%% These are more robust when using  xelatex  but may be employed with  pdflatex
%% The following definitons are meant to be re-defined in a style with \renewcommand
\newcommand{\divisionfont}{\relax}
\newcommand{\blocktitlefont}{\relax}
\newcommand{\contentsfont}{\relax}
\newcommand{\pagefont}{\relax}
\newcommand{\tabularfont}{\relax}
\newcommand{\xreffont}{\relax}
\newcommand{\titlepagefont}{\relax}
%% begin: font information supplied by "font-pdflatex-style" template
\usepackage{lmodern}
\usepackage[T1]{fontenc}
%% begin: font information supplied by "font-pdflatex-style" template
%% end: font setup and configuration for use with pdflatex
%% end: pdflatex-specific configuration
}
%% Symbols, align environment, commutative diagrams, bracket-matrix
\usepackage{amsmath}
\usepackage{amscd}
\usepackage{amssymb}
%% allow page breaks within display mathematics anywhere
%% level 4 is maximally permissive
%% this is exactly the opposite of AMSmath package philosophy
%% there are per-display, and per-equation options to control this
%% split, aligned, gathered, and alignedat are not affected
\allowdisplaybreaks[4]
%% allow more columns to a matrix
%% can make this even bigger by overriding with  latex.preamble.late  processing option
\setcounter{MaxMatrixCols}{30}
%%
%%
%% Division Titles, and Page Headers/Footers
%% titlesec package, loading "titleps" package cooperatively
%% See code comments about the necessity and purpose of "explicit" option.
%% The "newparttoc" option causes a consistent entry for parts in the ToC 
%% file, but it is only effective if there is a \titleformat for \part.
%% "pagestyles" loads the  titleps  package cooperatively.
\usepackage[explicit, newparttoc, pagestyles]{titlesec}
%% The companion titletoc package for the ToC.
\usepackage{titletoc}
%% Fixes a bug with transition from chapters to appendices in a "book"
%% See generating XSL code for more details about necessity
\newtitlemark{\chaptertitlename}
%% begin: customizations of page styles via the modal "titleps-style" template
%% Designed to use commands from the LaTeX "titleps" package
%% Plain pages should have the same font for page numbers
\renewpagestyle{plain}{%
\setfoot{}{\pagefont\thepage}{}%
}%
%% Single pages as in default LaTeX
\renewpagestyle{headings}{%
\sethead{\pagefont\slshape\MakeUppercase{\ifthechapter{\chaptertitlename\space\thechapter.\space}{}\chaptertitle}}{}{\pagefont\thepage}%
}%
\pagestyle{headings}
%% end: customizations of page styles via the modal "titleps-style" template
%%
%% Create globally-available macros to be provided for style writers
%% These are redefined for each occurence of each division
\newcommand{\divisionnameptx}{\relax}%
\newcommand{\titleptx}{\relax}%
\newcommand{\subtitleptx}{\relax}%
\newcommand{\shortitleptx}{\relax}%
\newcommand{\authorsptx}{\relax}%
\newcommand{\epigraphptx}{\relax}%
%% Create environments for possible occurences of each division
%% Environment for a PTX "chapter" at the level of a LaTeX "chapter"
\NewDocumentEnvironment{chapterptx}{mmmmmm}
{%
\renewcommand{\divisionnameptx}{Chapter}%
\renewcommand{\titleptx}{#1}%
\renewcommand{\subtitleptx}{#2}%
\renewcommand{\shortitleptx}{#3}%
\renewcommand{\authorsptx}{#4}%
\renewcommand{\epigraphptx}{#5}%
\chapter[{#3}]{#1}%
\label{#6}%
}{}%
%% Environment for a PTX "section" at the level of a LaTeX "section"
\NewDocumentEnvironment{sectionptx}{mmmmmm}
{%
\renewcommand{\divisionnameptx}{Section}%
\renewcommand{\titleptx}{#1}%
\renewcommand{\subtitleptx}{#2}%
\renewcommand{\shortitleptx}{#3}%
\renewcommand{\authorsptx}{#4}%
\renewcommand{\epigraphptx}{#5}%
\section[{#3}]{#1}%
\label{#6}%
}{}%
%%
%% Styles for six traditional LaTeX divisions
\titleformat{\part}[display]
{\divisionfont\Huge\bfseries\centering}{\divisionnameptx\space\thepart}{30pt}{\Huge#1}
[{\Large\centering\authorsptx}]
\titleformat{\chapter}[display]
{\divisionfont\huge\bfseries}{\divisionnameptx\space\thechapter}{20pt}{\Huge#1}
[{\Large\authorsptx}]
\titleformat{name=\chapter,numberless}[display]
{\divisionfont\huge\bfseries}{}{0pt}{#1}
[{\Large\authorsptx}]
\titlespacing*{\chapter}{0pt}{50pt}{40pt}
\titleformat{\section}[hang]
{\divisionfont\Large\bfseries}{\thesection}{1ex}{#1}
[{\large\authorsptx}]
\titleformat{name=\section,numberless}[block]
{\divisionfont\Large\bfseries}{}{0pt}{#1}
[{\large\authorsptx}]
\titlespacing*{\section}{0pt}{3.5ex plus 1ex minus .2ex}{2.3ex plus .2ex}
\titleformat{\subsection}[hang]
{\divisionfont\large\bfseries}{\thesubsection}{1ex}{#1}
[{\normalsize\authorsptx}]
\titleformat{name=\subsection,numberless}[block]
{\divisionfont\large\bfseries}{}{0pt}{#1}
[{\normalsize\authorsptx}]
\titlespacing*{\subsection}{0pt}{3.25ex plus 1ex minus .2ex}{1.5ex plus .2ex}
\titleformat{\subsubsection}[hang]
{\divisionfont\normalsize\bfseries}{\thesubsubsection}{1em}{#1}
[{\small\authorsptx}]
\titleformat{name=\subsubsection,numberless}[block]
{\divisionfont\normalsize\bfseries}{}{0pt}{#1}
[{\normalsize\authorsptx}]
\titlespacing*{\subsubsection}{0pt}{3.25ex plus 1ex minus .2ex}{1.5ex plus .2ex}
\titleformat{\paragraph}[hang]
{\divisionfont\normalsize\bfseries}{\theparagraph}{1em}{#1}
[{\small\authorsptx}]
\titleformat{name=\paragraph,numberless}[block]
{\divisionfont\normalsize\bfseries}{}{0pt}{#1}
[{\normalsize\authorsptx}]
\titlespacing*{\paragraph}{0pt}{3.25ex plus 1ex minus .2ex}{1.5em}
%%
%% Styles for five traditional LaTeX divisions
\titlecontents{part}%
[0pt]{\contentsmargin{0em}\addvspace{1pc}\contentsfont\bfseries}%
{\Large\thecontentslabel\enspace}{\Large}%
{}%
[\addvspace{.5pc}]%
\titlecontents{chapter}%
[0pt]{\contentsmargin{0em}\addvspace{1pc}\contentsfont\bfseries}%
{\large\thecontentslabel\enspace}{\large}%
{\hfill\bfseries\thecontentspage}%
[\addvspace{.5pc}]%
\dottedcontents{section}[3.8em]{\contentsfont}{2.3em}{1pc}%
\dottedcontents{subsection}[6.1em]{\contentsfont}{3.2em}{1pc}%
\dottedcontents{subsubsection}[9.3em]{\contentsfont}{4.3em}{1pc}%
%%
%% Begin: Semantic Macros
%% To preserve meaning in a LaTeX file
%%
%% \mono macro for content of "c", "cd", "tag", etc elements
%% Also used automatically in other constructions
%% Simply an alias for \texttt
%% Always defined, even if there is no need, or if a specific tt font is not loaded
\newcommand{\mono}[1]{\texttt{#1}}
%%
%% Following semantic macros are only defined here if their
%% use is required only in this specific document
%%
%% End: Semantic Macros
%% Division Numbering: Chapters, Sections, Subsections, etc
%% Division numbers may be turned off at some level ("depth")
%% A section *always* has depth 1, contrary to us counting from the document root
%% The latex default is 3.  If a larger number is present here, then
%% removing this command may make some cross-references ambiguous
%% The precursor variable $numbering-maxlevel is checked for consistency in the common XSL file
\setcounter{secnumdepth}{3}
%%
%% AMS "proof" environment is no longer used, but we leave previously
%% implemented \qedhere in place, should the LaTeX be recycled
\newcommand{\qedhere}{\relax}
%%
%% A faux tcolorbox whose only purpose is to provide common numbering
%% facilities for most blocks (possibly not projects, 2D displays)
%% Controlled by  numbering.theorems.level  processing parameter
\newtcolorbox[auto counter, number within=section]{block}{}
%%
%% This document is set to number PROJECT-LIKE on a separate numbering scheme
%% So, a faux tcolorbox whose only purpose is to provide this numbering
%% Controlled by  numbering.projects.level  processing parameter
\newtcolorbox[auto counter, number within=section]{project-distinct}{}
%% A faux tcolorbox whose only purpose is to provide common numbering
%% facilities for 2D displays which are subnumbered as part of a "sidebyside"
\newtcolorbox[auto counter, number within=tcb@cnt@block, number freestyle={\noexpand\thetcb@cnt@block(\noexpand\alph{\tcbcounter})}]{subdisplay}{}
%% Localize LaTeX supplied names (possibly none)
\renewcommand*{\chaptername}{Chapter}
%% Equation Numbering
%% Controlled by  numbering.equations.level  processing parameter
%% No adjustment here implies document-wide numbering
\numberwithin{equation}{section}
%% "tcolorbox" environment for a single image, occupying entire \linewidth
%% arguments are left-margin, width, right-margin, as multiples of
%% \linewidth, and are guaranteed to be positive and sum to 1.0
\tcbset{ imagestyle/.style={bwminimalstyle} }
\NewTColorBox{image}{mmm}{imagestyle,left skip=#1\linewidth,width=#2\linewidth}
%% More flexible list management, esp. for references
%% But also for specifying labels (i.e. custom order) on nested lists
\usepackage{enumitem}
%% hyperref driver does not need to be specified, it will be detected
%% Footnote marks in tcolorbox have broken linking under
%% hyperref, so it is necessary to turn off all linking
%% It *must* be given as a package option, not with \hypersetup
\usepackage[hyperfootnotes=false]{hyperref}
%% Hyperlinking active in electronic PDFs, all links solid and blue
\hypersetup{colorlinks=true,linkcolor=blue,citecolor=blue,filecolor=blue,urlcolor=blue}
\hypersetup{pdftitle={Gauss Contest : Problems and Solutions}}
%% If you manually remove hyperref, leave in this next command
\providecommand\phantomsection{}
%% If tikz has been loaded, replace ampersand with \amp macro
%% extpfeil package for certain extensible arrows,
%% as also provided by MathJax extension of the same name
%% NB: this package loads mtools, which loads calc, which redefines
%%     \setlength, so it can be removed if it seems to be in the 
%%     way and your math does not use:
%%     
%%     \xtwoheadrightarrow, \xtwoheadleftarrow, \xmapsto, \xlongequal, \xtofrom
%%     
%%     we have had to be extra careful with variable thickness
%%     lines in tables, and so also load this package late
\usepackage{extpfeil}
%% Custom Preamble Entries, late (use latex.preamble.late)
%% Begin: Author-provided packages
%% (From  docinfo/latex-preamble/package  elements)
%% End: Author-provided packages
%% Begin: Author-provided macros
%% (From  docinfo/macros  element)
%% Plus three from MBX for XML characters

\newcommand{\lt}{<}
\newcommand{\gt}{>}
\newcommand{\amp}{&}
%% End: Author-provided macros
\begin{document}
This contains a collection of problems that appeared in the National Mathematics Talent Contests (NMTC), a mathematics competition conducted for students in India. There are different levels depending on the age of students. For students in the age group 10-11, the competition is called Gauss contest. The problems that appeared in Gauss contest during the years 2017-2019 are solved and published here.%
%
%
\typeout{************************************************}
\typeout{Chapter 1 2019}
\typeout{************************************************}
%
\begin{chapterptx}{2019}{}{2019}{}{}{x:chapter:problems-2019}
%
%
\typeout{************************************************}
\typeout{Section 1.1 Problem-1}
\typeout{************************************************}
%
\begin{sectionptx}{Problem-1}{}{Problem-1}{}{}{x:section:gauss_2019_problem_1}
%
\begin{enumerate}
\item{}Find positive integers \(m, n\) such that:%
\begin{equation*}
\cfrac{1}{m} + \cfrac{1}{n} = \cfrac{3}{17}
\end{equation*}
%
\item{}Find positive integers \(m, n, p\) such that:%
\begin{equation*}
\cfrac{1}{m} + \cfrac{1}{n} + \cfrac{1}{p} = \cfrac{3}{17}
\end{equation*}
%
\end{enumerate}
\textbf{\blocktitlefont Hint}.\quad{}Can \(m\) and \(n\) both be multiples of \(17\)?%
\textbf{\blocktitlefont Solution}.\quad{}%
\begin{enumerate}
\item{}If \(m = n = 17\), we get:%
\begin{equation*}
\cfrac{1}{m} + \cfrac{1}{n} = \cfrac{2}{17}
\end{equation*}
If both \(m\) and \(n\) are multiples of \(17\), say \(m = 17a,\ n = 17b\), such that \(a \geq 1,\ b \geq 1\), then we have:%
\begin{equation*}
\cfrac{1}{m} + \cfrac{1}{n} = \cfrac{1}{17} \Bigg( \cfrac{1}{a} + \cfrac{1}{b} \Bigg) \leq \cfrac{2}{17}
\end{equation*}
Thus we see that the both \(m\) and \(n\) cannot be multiples of \(17\) if they are to solve this equation. But it is also clear that at least one of them must be a multiple of \(17\). So let \(n = 17k\). Making this substitution, we get:%
\begin{align*}
\cfrac{1}{m} + \cfrac{1}{n} \amp= \cfrac{3}{17}\\
\cfrac{1}{m} + \cfrac{1}{17k} \amp= \cfrac{3}{17}\\
\cfrac{1}{m} \amp= \cfrac{3}{17} - \cfrac{1}{17k}\\
\amp= \cfrac{3k - 1}{17k}\\
\implies m \amp= \cfrac{17k}{3k - 1}
\end{align*}
From this we can conclude that \(3k - 1\) must divide \(17k\). One obvious choice is for \(3k - 1 = 17\) which gives \(k = 6\). So a solution for \((m, n)\) is:%
\begin{equation*}
(m, n) = (6, 6 \times 17)
\end{equation*}
%
\item{}If \(m = n = p = 17\), we get:%
\begin{equation*}
\cfrac{1}{m} + \cfrac{1}{n} + \cfrac{1}{p} = \cfrac{3}{17}
\end{equation*}
So, this is one obvious solution. But there are other interesting solutions as well. For finding them, we don't need to solve the system all over again. We can piggyback on the solution of the first part. First, let us write down the solution to the first part:%
\begin{equation*}
\cfrac{1}{6} + \cfrac{1}{6 \times 17} = \cfrac{3}{17}
\end{equation*}
We now make use of the simple trick: multiply and divide the second term on the LHS by \(3\). It becomes:%
\begin{align*}
\cfrac{1}{6} + \cfrac{1}{6 \times 17} \times \cfrac{3}{3} \amp= \cfrac{3}{17}\\
\cfrac{1}{6} + \cfrac{1}{6 \times 3} \times \cfrac{3}{17} \amp= \cfrac{3}{17}
\end{align*}
Now, we have a \(\frac{3}{17}\) appear on the LHS again. We just substitute it with the solution from the first part:%
\begin{align*}
\cfrac{1}{6} + \cfrac{1}{6 \times 3} \times \cfrac{3}{17} \amp= \cfrac{3}{17}\\
\cfrac{1}{6} + \cfrac{1}{6 \times 3} \times \Bigg( \cfrac{1}{6} + \cfrac{1}{6 \times 17} \Bigg) \amp= \cfrac{3}{17}\\
\cfrac{1}{6} + \cfrac{1}{6^2 \times 3} + \cfrac{1}{6^2 \times 3 \times 17}  \amp= \cfrac{3}{17}
\end{align*}
From this we can easily read the solution in three variables:%
\begin{equation*}
m = 6,\ n = 3 \times 6^2,\ p = 3 \times 6^2 \times 17
\end{equation*}
%
\end{enumerate}
%
\end{sectionptx}
%
%
\typeout{************************************************}
\typeout{Section 1.2 Problem-2}
\typeout{************************************************}
%
\begin{sectionptx}{Problem-2}{}{Problem-2}{}{}{x:section:gauss_2019_problem_2}
Find the largest positive integer \(n\) such that \(3^{n}\) divides the \(999\) digit number \(999\cdots99\).%
\textbf{\blocktitlefont Hint}.\quad{}Dividing the number by \(9\) is a good start.%
\textbf{\blocktitlefont Solution}.\quad{}We can see right away that the \(999\) digit number will be divisible by \(9\). So let us go ahead and divide it by \(9\). This will give us the \(999\) digit number where all digits are \(1\). At this stage, we have the factorization:%
\begin{equation*}
\underbrace{99 \cdots 99}_{999 \text{ nines}} = 9 \times \underbrace{11 \cdots 11}_{999 \text{ ones}}
\end{equation*}
We know that \(\underbrace{11 \cdots 11}_{999 \text{ ones}}\) is also divisible by \(9\) as the sum of its digits is \(999\) which is divisible by \(9\). So we can go ahead and divide it by \(9\). To keep things more tractable, let us begin by dividing \(\underbrace{11 \cdots 11}_{9 \text{ ones}}\) by \(9\):%
\begin{equation*}
\cfrac{111111111}{9} = 12345679
\end{equation*}
Now, if we move to the \textdollar{}18\textdollar{} digit number:%
\begin{equation*}
\cfrac{\underbrace{11 \cdots 11}_{18 \text{ ones}}}{9} = 012345679\ 012345679
\end{equation*}
Note that we have attached a zero at the beginning for the sake of maintaining symmetry. With this, we can quickly generalize this to the \(999\) digit number:%
\begin{equation*}
\cfrac{\underbrace{11 \cdots 11}_{999 \text{ ones}}}{9} = \underbrace{012345679\ \cdots\ 012345679}_{111 \text{ blocks of } 012345679}
\end{equation*}
Every set of \(9\) contiguous ones contribute to one block of \(012345679\) upon division. There are \(111\) such blocks. Now, let us see if this latest quotient we have obtained is divisible by \(3\). For that, we need to compute the digit sum:%
\begin{equation*}
(0 + 1 + 2 + 3 + 4 + 5 + 6 + 7 + 9) \times 111 = 37 \times 111 = 3 \times 37^2
\end{equation*}
We see that it is divisible by \(3\) but not by \(9\) and therefore not divisible by any power of \(3\) greater than or equal to \(2\). So we can write:%
\begin{equation*}
\underbrace{012345679\ \cdots\ 012345679}_{111 \text{ blocks of } 012345679} = 3k
\end{equation*}
Here \(k\) is some positive integer that is not a multiple of \(3\). With that, we have:%
\begin{align*}
\underbrace{99 \cdots 99}_{999 \text{ nines}} \amp= 9 \times \underbrace{11 \cdots 11}_{999 \text{ ones}}\\
\amp= 9 \times 9 \times \underbrace{012345679\ \cdots\ 012345679}_{111 \text{ blocks of } 012345679}\\
\amp= 9 \times 9 \times 3k\\
\amp= 3^{2 + 2 + 1}k\\
\amp= 3^5k
\end{align*}
Therefore, the largest positive integer \(n\) such that \(3^n\) divides the given \(999\) digit number is \(n = 5\).%
\end{sectionptx}
%
%
\typeout{************************************************}
\typeout{Section 1.3 Problem-3}
\typeout{************************************************}
%
\begin{sectionptx}{Problem-3}{}{Problem-3}{}{}{x:section:gauss_2019_problem_3}
Inside a square of area \(36\ cm^2\), there are shaded regions as shown. \begin{image}{0}{1}{0}%
\includegraphics[width=\linewidth]{../images/i1.png}
\end{image}%
 The ratio of the shaded area to the unshaded area is \(3:1\). What is the value of \(a + b + c + d\) where \(a, b, c, d\) are the lengths of the bases of the shaded regions? Further, if three of \(a, b, c, d\) are equal integers and one different, then find them.%
\textbf{\blocktitlefont Hint}.\quad{}Dividing the number by \(9\) is a good start.%
\textbf{\blocktitlefont Solution}.\quad{}We obtain four triangles from the shaded area. The sum of the area of these shaded triangles is:%
\begin{equation*}
\cfrac{1}{2} \cdot a \cdot 6 + \Bigg( \cfrac{1}{2} \cdot b \cdot 6 + \cfrac{1}{2} \cdot c \cdot 6 \Bigg) + \cfrac{1}{2} \cdot d \cdot 6 = \cfrac{1}{2} \cdot 6 \cdot (a + b + c + d)
\end{equation*}
As the ratio of the shaded to unshaded area is \(3:1\), the area of the shaded region will be \(\frac{3}{4} \times 36 = 27\ cm^2\). With this, we get:%
\begin{equation*}
\cfrac{1}{2} \cdot 6 \cdot (a + b + c + d) = 27 \implies a + b + c + d = 9
\end{equation*}
As for the second part of the question, some three integers are equal. Let the four numbers be \((x, x, x, y)\). Then, we have:%
\begin{equation*}
3x + y = 9 \implies y = 9 - 3x
\end{equation*}
\(x\) can only take values \(1\) or \(2\) as we want \(y > 0\). For \(x = 1\), we get \(y = 6\). But the side of the square is itself of length \(6\). So this solution is ruled out. For \(x = 2\), we get \(y = 3\). So the four sides are \((2, 2, 2, 3)\). Note that the order doesn't matter.%
\end{sectionptx}
\end{chapterptx}
\end{document}